\documentclass[10pt,a4paper]{amsart}
\usepackage[margin=19mm]{geometry}
\usepackage{amsmath,amssymb,mathtools}
\usepackage[scr=rsfs,cal=euler]{mathalfa}
\usepackage{xfrac}
\usepackage[unicode,hidelinks,bookmarks=false]{hyperref}
\newcommand{\ie}{i.e.\ }
\newcommand{\cf}{cf.\ }
\newcommand{\resp}{respectively\ }
\newcommand{\Subsection}{\subsection}
\providecommand{\nobreakdash}{\nobreak\hbox{-}}
\setlength{\emergencystretch}{2em}
\multlinegap0pt
\usepackage{bm}
\usepackage{bbm}
\usepackage{dsfont}
\usepackage{xspace}
\usepackage{cancel}
\usepackage{mathtools}
\usepackage{amsthm}
\makeatletter
\@ifundefined{thm}{\newtheorem{thm}{Theorem}}{}
\makeatother
\newcommand{\cl}{\operatorname{cl}}
\renewcommand{\int}{\operatorname{int}}
\title[Homeomorphisms between compact subsets of real numbers]{Homeomorphisms between compact subsets of real numbers}
\author{S{\l}awomir Kusi\'nski, Szymon Plewik}
\date{Source: arXiv:2602.15542v2 (CC BY 4.0)}
\begin{document}
\maketitle

\noindent\emph{Source and licence.} Homeomorphisms between compact subsets of real numbers. Version arXiv:2602.15542v2 (CC BY 4.0), distributed under the Creative Commons Attribution 4.0 International licence, as recorded in the arXiv licence metadata for this exact version and shown on the abstract page https://arxiv.org/abs/2602.15542v2. Full rights notice: licenses/arxiv-2602.15542-CC-BY-4.0.txt.\par\smallskip

\noindent\textit{Editorial scope.} Counted unit: the Thm whose source label is \texttt{tr1} in the article (source0.tex lines 135--138, source label \texttt{tr1}), delivered together with the article's own complete original proof of it (source0.tex lines 139--160, one \texttt{proof} environment of 22 lines). No mathematics is written, repaired or re-derived: statement and proof are the article's own text line for line. Required context retained uncounted: none. Every definition, display and label of those lines is retained unchanged. Omitted from this package and not redistributed: 1--134, 161--587. Nothing inside a retained excerpt is omitted or altered: every retained body is delivered exactly as the article's lines are written, so no omission receipt of any kind accompanies this package. Citations: the retained excerpts carry no citation command at all, so no bibliography entry of the article is transcribed and no reference is redistributed. Reference adaptation: none was needed and none was made; every reference inside the delivered unit resolves inside it, so no unresolved reference or citation is produced. Printed slips kept literal: no printed word, symbol, hypothesis or number of the counted statement or of its proof was altered. Layout and class: the article's own class is \texttt{amsart} with packages including inputenc, amsfonts, amsmath, amssymb, amsthm, indentfirst, xcolor, fontenc; the wrapper is the lane's pinned \texttt{amsart} editorial wrapper at 10pt on A4, which restates the theorem-like environments the retained text needs and adds the article's own macro definitions that the retained text needs (\texttt{\textbackslash cl}, \texttt{\textbackslash int}), with extra wrapper packages (bm, bbm, dsfont, xspace, cancel, mathtools). The delivered numbering is the unit's own. Assets: no retained span contains a figure environment, a graphics-inclusion command, a PDF-page inclusion, a table or a TikZ picture; no image file is copied into this package, the package declares no asset dependency and no image file is redistributed with it. Licence: version arXiv:2602.15542v2 of this article is distributed under the Creative Commons Attribution 4.0 International licence, as recorded in the arXiv licence metadata for this exact version and shown on the abstract page https://arxiv.org/abs/2602.15542v2. A full rights notice is delivered at licenses/arxiv-2602.15542-CC-BY-4.0.txt. Characters: every retained excerpt is pure ASCII, so the delivered engine, XeLaTeX with its default Latin Modern text font, reports no missing character.
   \begin{thm}\label{tr1}
 Let $(\mathcal  G \cup \mathcal I, <)$ be a countable  ordered space such that  the sets   $ \mathcal G$ and  $\mathcal I $ are disjoint. If $ \mathcal G$ interlaces  $\mathcal I $, and  the smallest and the largest element of $\mathcal  G \cup \mathcal I$ belong to $\mathcal G $, then there exist a compact subset $Y\subset \mathbb R$ and the isomorphism  $F\colon \mathcal G\cup \mathcal I \to \mathcal Q_Y$ such that  $F[\mathcal G] =\mathcal G_Y$ and 
$ F[\mathcal I] =\mathcal I_Y$.  
   \end{thm}

   \begin{proof}  Give $(\mathcal G\cup\mathcal I)\times\mathbb Q$ the lexicographic order, where $\mathbb Q$ denotes the set of all rational numbers. Clearly, this order is countable, dense and has no endpoints, hence its Dedekind completion $\mathfrak C$ is isomorphic to $[0,1]$. Let $H\colon \mathfrak C \to [0,1]$ be such an isomorphism.
For each $I\in \mathcal I$ (resp. $G\in \mathcal G$), put  
$$
B_I=\cl_{\mathfrak C}(\{I\}\times\mathbb Q) \quad \mbox{ (resp. } B_G=\int_{\mathfrak C} \cl_\mathfrak C(\{G\}\times\mathbb Q))$$ and 
define
$$\textstyle 
Y=H\left[\mathfrak C\setminus\bigcup_{G\in\mathcal G}B_G\right].$$
   \indent  If $I\in\mathcal I$, then $H[B_I]$, being a closed interval,  is a non-trivial component of $Y$. Indeed,  if a connected subset of $Y$ met both 
	$H[B_I]$ and another $H[B_{J}]$, the interlacing hypothesis would place some $G\in\mathcal G$ between $I$ and $J$, and the removed open interval $H[B_G]$ would separate the two parts.  Thus, any connected subset of $Y$ meets at most one $H[B_I]$, where $I\in\mathcal I$.	
	Moreover, every non-trivial interval $[a,b] \subseteq Y$ has to be contained in  $H[B_I]$ for some $I\in\mathcal I$, since
	$\mathcal G$ is countable and $H^{-1}([a,b])$ has the cardinality continuum. 
	
	
			If $G\in\mathcal G$, then $H[B_G]$ is a  component of $[0,1] \setminus Y$, since $H[B_G]$ is an interval with endpoints in $Y\cup \{0,1\}$.
			
			Define the necessary isomorphism   $F: \mathcal G\cup \mathcal I \to \mathcal Q_Y$  as follows.
  \[F(J)= \begin{cases}
	  (\leftarrow, \min Y),& \mbox{ for } J= \min \mathcal G,\\
		( \max Y,\rightarrow), & \mbox{ for } J=\max  \mathcal G,\\
		H[B_J],& \mbox{ for any other } J\in \mathcal G \cup \mathcal  I.
	\end{cases}\]			 
   \end{proof} 

\end{document}
